\section{Method}
\label{sec:method}

The controller is a two-level cascade. A centralised planner plans the motion of the ring and of every link over a
horizon; a tracking MPC on each carrier follows that carrier's share of the plan. A change of team is handled
inside the planner, by changing the problem it solves.

\subsection{Cooperative planner and trackers}
\label{sec:planner}

The planner's model is the ring as a rigid body, pose $(\mathbf{p}_L, \mathbf{q}_L)$ and twist $(\mathbf{v}_L,
\boldsymbol{\omega}_L)$, with each of the $N$ carriers reduced to a link of length $l_i$ attached at the plate
point $\boldsymbol{\rho}_i$ in the ring frame, with unit direction $\mathbf{q}_i$ and tension $\lambda_i$. The
carrier positions follow kinematically, $\mathbf{p}_i = \mathbf{p}_L + R(\mathbf{q}_L)\boldsymbol{\rho}_i -
l_i\mathbf{q}_i$, and the thrust each must produce follows from its link's acceleration and tension. At every
tick the planner solves
\begin{equation}
\begin{aligned}
  \min_{\mathbf{x}_{0:N_h},\,\mathbf{u}_{0:N_h-1}} \;& \sum_{k=0}^{N_h-1} \|\mathbf{y}_k - \mathbf{y}_k^{\mathrm{ref}}\|_W^2
    + \|\mathbf{y}_{N_h} - \mathbf{y}_{N_h}^{\mathrm{ref}}\|_{W_{N_h}}^2 \\
  \text{s.t.}\;& \mathbf{x}_{k+1} = \mathbf{F}(\mathbf{x}_k, \mathbf{u}_k;\, \{\boldsymbol{\rho}_i, l_i\}, \mathbf{d}),
    \quad C_0\mathbf{x}_0 = C_0\hat{\mathbf{x}}(t), \\
  & \lambda_{\min} \le \lambda_{i,k} \le \lambda_{\max}, \quad f_{\min} \le \|\mathbf{f}_{i,k}\| \le f_{\max},
\end{aligned}
\label{eq:ocp}
\end{equation}
where the output $\mathbf{y}$ stacks the ring's position, velocity, attitude error and rate, the link directions,
tensions and rates, and the inputs; $\mathbf{f}_{i,k}$ is carrier $i$'s planned thrust, its bounds softened with
slack. The horizon has $N_h = 20$ nodes over 2\,s, so a node is the 100\,ms planner tick. The constraint on node~0
pins only what motion capture measures, the ring's pose and twist and the link directions; the tensions and link
rates are left free and warm-started from the previous solution. A lumped force $\mathbf{d}$ on the ring, formed
from the integral of its position error in steady flight, removes steady offsets through the model
\cite{pannocchia2003disturbance, maeder2009linear}. The references follow from the ring's trajectory by flatness,
with the nominal formation tilted along the effective gravity and the tensions scaled with it.

Problem~\eqref{eq:ocp} is solved with acados \cite{verschueren2022acados} in a real-time iteration scheme
\cite{diehl2002rti}: a fixed number of Gauss--Newton SQP iterations per tick (one unless stated;
Section~\ref{sec:res_scale}), warm-started, with HPIPM \cite{frison2020hpipm}
solving the QP. On the rig no new iteration starts after 60\,ms of the tick, and a plan not ready by then is replaced
by the previous one shifted by a node.

Each carrier's tracker is an MPC on a 17-state model in which the four stick commands are states and their rates the
inputs, with the flight controller's rate loop modelled as first order and the identified affine thrust map. It
solves every 20\,ms over the same 2\,s horizon, following the planner's position, velocity and thrust direction. The
link's pull enters its model as a known parameter along the horizon: the planned pull, shifted by the difference
between the pull measured from motion capture and the planned one at the first node, so that the tracker anticipates
the link's force rather than reacting to it.

\subsection{Reconfiguration: release and drop}
\label{sec:reconfig}

Two properties of \eqref{eq:ocp} make an in-flight change of team possible. The plate points
$\boldsymbol{\rho}_i$ and link lengths $l_i$ are runtime parameters, so only the team size fixes the problem's
dimensions, and one solver per size, built in advance and primed on the measured state, covers any layout. And the
survivors keep their own plates: the nominal tensions come from a moment balance at the actual plate points, so the
ring's reference stays level without re-spacing the team.

\emph{Detection.} A release the planner is told of is a DETACH command, which switches the leaver's magnet off and
resizes on the same tick. An unannounced release is detected from the link alone. A rigid link on its plate cannot
measure longer than itself, so when the leaver's link, measured by motion capture between the carrier's pivot and
its plate point, reads more than 6\,cm over its length for two consecutive ticks, the planner removes that carrier.
No force sensor and no message from the leaver are needed.

\emph{Resize.} At the event the planner swaps to the solver for $N-1$ carriers, sets the survivors' plate points and
link lengths, and warm-starts it from the measured state with the tensions and link directions free at node~0. The
first solve redistributes the load; there is no separate hand-off phase. With three survivors the vertical shares
are fixed by statics: each is the ring's weight times the barycentric coordinate of its centre in the survivors'
triangle, and a level ring exists only while the centre lies inside the polygon of the plates
\cite{michael2011cooperative}. For plates on a circle of radius $r$, with $g_{\max}$ the survivors' largest angular
gap, the support margin
\begin{equation}
  m = r\cos(g_{\max}/2)
  \label{eq:margin}
\end{equation}
is the distance from the centre to the nearest edge. Losing one of four evenly spaced carriers puts the centre on an
edge, so the four-carrier layout uses plates 3, 7, 9 and 11, where losing plate 9 leaves an even triangle with a
112\,mm margin.

\emph{Drop rule.} A loss that leaves fewer survivors than the configured minimum (three by default), or a ring tilt
above $30^\circ$, drops the ring: every magnet is switched off and every carrier lands. Release chains that continue
below three carriers fly on while the survivors' largest gap is at most $200^\circ$, which keeps a T and an opposite
pair. The freed carrier steps out radially, clear of the ring, and lands and disarms alone while the survivors carry
on.

\subsection{Rejoin of a departed carrier}
\label{sec:rejoin}

A departed carrier can return to its own plate. At its release the planner records its plate point and the link's
measured length. \emph{Staging:} the carrier is held on its plate's bearing 0.3\,m beyond the formation's reach,
the point moving and turning with the planned ring and kept inside the position limit. \emph{Approach:} on the
rejoin command, given only while the ring holds its path, the carrier descends at 0.1\,m/s so that its link's
lower end lands on the plate point of the measured ring, with a feedback loop on the measured end correcting the
end's height. \emph{Weld:} the magnet is switched on once the end has stayed within 1.5\,cm of its plate point on
every sample over 0.3\,s (at least eight samples) at under 2\,cm/s; the joint must then report attached, or the
carrier is let go and backs out to staging. \emph{Fold-in:} the planner grows to $N+1$ with the plate point and the
link re-measured from pivot to plate; the rejoiner is seeded at 0.1\,N and every tension is blended to the balanced
split over 10\,s, with its link feedforward ramped over 1.5\,s and the ring's path held throughout, as a dock blends
thrust across a switch \cite{sugihara2023trady}. Through the blend the rejoiner's height reference is capped at the
weld height, so that it cannot lift its plate and tilt the ring before it carries its share, and it counts as absent
for the drop rule. LAND or a drop sends every staged or rejoining carrier up, out along its bearing and down alone.

\subsection{Attach of a newcomer}
\label{sec:attach}

A newcomer is flown by another team's controller down to contact; its magnet catches only there. Our side admits it
from motion capture alone, with its magnet tip expressed in the plate's frame. A \emph{contact} is the tip within
24\,mm of the plate point and close to the plate's face, nearly still for 0.3\,s. A contact is only a candidate, since
a tip resting on a plate and a caught magnet look alike, so the catch is \emph{confirmed} by co-motion: after the
newcomer's small lift, the tip has stayed within 6\,mm of its contact point while the plate moved with it, on a taut
tether. Only then do our trackers take the newcomer over and the planner grows to $N+1$ with the recorded plate
point, the newcomer seeded at 0.1\,N and blended as in Section~\ref{sec:rejoin}. A newcomer not confirmed in time is
taken back: our tracker latches onto its measured state, lifts it clear and lands it alone.
